\title{Six Birds Foundations II:\\[4pt]
Admissibility, Meta-Theory, and the Exact-Six Program}
\author{Ioannis Tsiokos\,\orcidlink{0009-0009-7659-5964}}
\date{Version 3: October 2, 2026\\[2pt] \footnotesize (Version 1: April 20, 2026)}

\maketitle
\thispagestyle{firstpage}

\begin{abstract}
\noindent
Six Birds Foundations I~\citep{Tsiokos2026SixBirdsFoundations}
described how a layer of description closes over what it tracks, and
identified six roles that such closure uses: descent (whether an update
respects a quotient), representability of a macro-level specification,
enriched route and coherence mismatch, staged order and refinement, packaging into
quotients, and audit. This paper gives the six roles precise meanings
and asks whether exactly six are needed. For each role it fixes what
counts as evidence for it, the level at which it acts, and where it
meets the other roles, and it states the bookkeeping rules a
description must follow before any role claim counts.

The question is posed on a precisely delimited domain, \FATCD{}, of finite
audited typed closure descriptions: finite collections of typed records in
which every claim carries an honest audit. Without further hypotheses the paper
proves three things. Each of the six roles occurs in some description of the
domain. Every description has a decomposition record with one channel per role,
each channel in one of five statuses (active, collapsed, below threshold,
outside scope, absent), and every feature is accounted for by an active
channel, an inactive channel-status entry, or one of eleven residual labels;
the last two, \emph{unresolved candidate} and \emph{recorded seventh-role
screen}, apply only when none of the first nine places the feature. The six channels are part of the record's definition, so the
count is of channels, not of active roles, and the record need not be unique.
Probability, causality and agency, the most natural candidates for a seventh
role, are placed without a seventh role in the forms the paper specifies.

The upper bound is conditional. Its recognition hypothesis is assumed, not
derived, and asks, without reference to any decomposition record, that every
feature be content of one of the six roles, an annotation, data that asserts
nothing, a checked restatement of such content, or content whose check is not
finite over the description (placed as outside-domain structure); the
upper-bound theorem shows that each of these ways yields one of the first nine
labels. Under the hypothesis a feature that no channel accounts for receives
one of labels (iii)--(ix), and no content record, annotation or leftover record
is licensed the last two. The hypothesis is a genuine restriction and depends
on presentation. A standalone formula stating that finitely many recorded
weights are nonnegative and sum to one fails it, while the same weights
recorded as the single row of a transition kernel, with a certificate that the
row is a probability distribution, satisfy it. On the other hand, any closed formula
of the paper's finite check language that does not mention the six witness
predicates can be recorded, provided the completed description lies in \FATCD{}, with its truth value, as a one-variable
specification over two records of the same kind, which the hypothesis accepts;
a false formula is recorded through its checked negation, and the rest of the
description must still satisfy the hypothesis. The paper does not claim six roles
for all of emergence, a full algebra of the six roles, their independence in
every setting, or an empirical realization.
Appendix~\ref{app:mechanization} describes the accompanying Lean
development and what it covers.
\end{abstract}
